\documentclass[11pt]{article}
\usepackage{preamble}
\setlength{\parskip}{4pt}

\begin{document}

\daybanner{AP Statistics \textperiodcentered\ Unit 2 \textperiodcentered\ Topic 2.3 \textperiodcentered\ B: 2026-10-20 \textperiodcentered\ E: 2026-11-04 \textperiodcentered\ TEACHER KEY}{What Did Each Simulation Estimate?}

\sectionbanner{CED TETHER}

{\small
\begin{itemize}[leftmargin=1.5em,itemsep=2pt,topsep=2pt]
  \item \textbf{Skill 3.A} --- Determine relative frequencies, proportions, or probabilities using simulation or calculations.
  \item \textbf{EU UNC-2} --- Simulation allows us to anticipate patterns in data.
  \item \textbf{LO UNC-2.A} --- Estimate probabilities using simulation. [Skill 3.A]
  \item \textbf{EK UNC-2.A.1} --- A random process generates results that are determined by chance.
  \item \textbf{EK UNC-2.A.2} --- An outcome is the result of a trial of a random process.
  \item \textbf{EK UNC-2.A.3} --- An event is a collection of outcomes.
  \item \textbf{EK UNC-2.A.4} --- Simulation is a way to model random events, such that simulated outcomes closely match real-world outcomes. All possible outcomes are associated with a value to be determined by chance. Record the counts of simulated outcomes and the count total.
  \item \textbf{EK UNC-2.A.5} --- The relative frequency of an outcome or event in simulated or empirical data can be used to estimate the probability of that outcome or event.
  \item \textbf{EK UNC-2.A.6} --- The law of large numbers states that simulated (empirical) probabilities tend to get closer to the true probability as the number of trials increases.
\end{itemize}
}
\textbf{Essential question:} How must chance labels, one trial, and the success event align with the probability question?\par
\textbf{DOK-3 skill:} 4.B \textperiodcentered\ \textbf{FRQ pattern:} {\small\texttt{whose-\allowbreak{}simulation-\allowbreak{}design-\allowbreak{}is-\allowbreak{}valid-\allowbreak{}and-\allowbreak{}what-\allowbreak{}it-\allowbreak{}misestimates}}\par
\textbf{DOK rationale:} Part (c) weighs competing claims using the day's numerical or graphical evidence and explains a limitation or consequence; the judgment requires linked reasoning beyond a calculation.\par

\frameworkphaseheader{Do Now (first take) $\rightarrow$ Explore (video + follow-along) $\rightarrow$ Exit (finish + turn in)}{1 $\rightarrow$ 3}{45}{%
\begin{itemize}[leftmargin=1.5em,itemsep=2pt,topsep=2pt]
  \item Show the board at the bell and collect a one-sentence first take without confirming a claim.
  \item Run the named video follow-along, then allow about 10 minutes for parts (a)--(c).
  \item Ask for evidence linked to a claim. Collect the sheet and hand-score part (c) with E/P/I.
\end{itemize}
}{%
\begin{itemize}[leftmargin=1.5em,itemsep=2pt,topsep=2pt]
  \item Choose a simulation result and name one design feature in a first take.
  \item Complete the follow-along about outcomes, events, chance assignments, trials, and relative frequency.
  \item Finish (a)--(c), revise the first take in the exit line, and turn in the sheet.
\end{itemize}
}{%
\begin{itemize}[leftmargin=1.5em,itemsep=2pt,topsep=2pt]
  \item Which number or visible feature supports your claim?
  \item Why does the other claim go beyond that evidence?
  \item What would you tell someone who chose the other claim?
\end{itemize}
}{Circulate and ask students to connect their choice to evidence. Score part (c) using the three required reasoning elements; do not require dropped extension work.}

\begin{callout}{\IconWarn\ WATCH FOR}{calloutred}
\begin{itemize}[leftmargin=1.5em,itemsep=2pt,topsep=2pt]
  \item A choice with copied numbers but no explanation of how they support it.
  \item Treating a model or average as a guaranteed individual outcome.
\end{itemize}
\end{callout}

\sectionbanner{SCORING PART (c) --- E / P / I}

\textbf{Expected elements}
\begin{itemize}[leftmargin=1.5em,itemsep=2pt,topsep=2pt]
  \item \textbf{reasoning-1} (required) --- Identifies the 0.10 versus 0.12 assignment mismatch
  \item \textbf{reasoning-2} (required) --- Identifies exactly 2 versus at least 2 and rejects choosing by the larger estimate
  \item \textbf{reasoning-3} (required) --- Gives five independent 00--99 pairs, 12 left-handed labels, and the at-least-2 event
\end{itemize}

{\small\renewcommand{\arraystretch}{1.25}
\noindent\begin{tabular}{|p{0.5in}|p{5.9in}|}
\hline
\textbf{E} & Includes all three required reasoning elements above, with correct contextual evidence. \\ \hline
\textbf{P} & Includes at least one substantive correct reasoning link above, but omits or misstates another required element. \\ \hline
\textbf{I} & Includes no substantive correct reasoning link above; an unsupported choice or copied number alone is insufficient. \\ \hline
\end{tabular}}

\textbf{Common mistakes}
\begin{itemize}[leftmargin=1.5em,itemsep=2pt,topsep=2pt]
  \item Counting digit 0 as 12 percent because 0 is one of the possible outcomes
  \item Treating exactly 2 and at least 2 as the same event
  \item Choosing Student 2 solely because two-digit numbers look more precise
  \item Fixing the label assignment but continuing to count exactly 2
\end{itemize}

\clearpage
\focusbanner{TODAY'S PROBLEM --- annotated}

Suppose each of five students can be treated as independently selected, and each has probability 0.12 of being left-handed. Two students use 1,000 simulation trials to estimate the probability that at least 2 of the 5 students are left-handed. Their plans and results are shown.\par
\begin{center}
\textbf{Results: Student 1 = 87/1,000; Student 2 = 94/1,000}\par\smallskip
\begin{tabular}{|l|c|c|}
\hline
\textbf{Plan} & \textbf{$L$ assignment} & \textbf{Trial; success} \\ \hline
\textbf{1} & 0 of 0--9 & 5 digits; $2+$ L \\ \hline
\textbf{2} & 00--11 of 00--99 & 5 pairs; exactly 2 L \\ \hline
\end{tabular}
\end{center}
\begin{firsttakebox}{FIRST TAKE (5 min)}
Which simulation result would you trust more for the question asked---Student 1, Student 2, or neither? Name one reason.\par
\answer{Not graded. Accept a tentative choice; the lesson should move students from decimal comparison to model-event alignment.}
\end{firsttakebox}

\textit{Rules callout ``THREE MATCHES MAKE A VALID SIMULATION (keep on the page)'' is printed on the student sheet.}\par\medskip

\dokbadge{1}~\textbf{(a)} List the left-handed counts included in at least 2 of 5. State the chance one student is left-handed.\par
\answer{The event includes 2, 3, 4, or 5 left-handed students. A valid random-label assignment must designate 12 percent of equally likely labels as left-handed.}

\dokbadge{2}~\textbf{(b)} Find each reported estimate. For each plan, identify either a mismatch in the chance assignment or a mismatch in the event counted.\par
\answer{Student 1 reports $87/1000=0.087$ and Student 2 reports $94/1000=0.094$. Student 1 assigns only $1/10=0.10$ of the digits to left-handed, although the success event is correctly at least 2. Student 2 assigns $12/100=0.12$ of the pairs to left-handed, but counts exactly 2 rather than 2, 3, 4, or 5.}

\dokbadge{3}~\textbf{(c)} Would you trust either plan for the question asked? Explain both mismatches and why the larger reported estimate does not validate a plan. Describe one corrected trial, including the random labels and the event counted. Write 3--4 sentences.\par
\answer{Neither plan matches the question. Student 1 assigns $1/10=0.10$, not $0.12$, to left-handed; Student 2 counts exactly 2 and leaves out 3, 4, and 5. A larger reported estimate does not fix either mismatch, and estimates also vary from run to run. A corrected trial draws five independent pairs from 00--99, uses 00--11 for left-handed and 12--99 for right-handed, and counts the trial when at least two selections are left-handed; repeat many trials and divide successes by trials.}

\begin{summaryexitbox}{EXIT}
One thing the video changed about my first take: \hrulefill
\end{summaryexitbox}

\end{document}
